% START: Metadaten ersetzen, dann die kommentierten Beispielabschnitte ändern.
% Bauen: im Repository-Hauptordner `make hausarbeit` (siehe README.md).
\documentclass[12pt,a4paper,ngerman]{scrartcl}
\usepackage[T1]{fontenc}
\usepackage[utf8]{inputenc}
\usepackage[ngerman]{babel}
\usepackage{lmodern}
\usepackage[a4paper,margin=2.5cm,headheight=18pt]{geometry}
\usepackage[bottom]{footmisc} % Fußnoten unten auf der zugehörigen Seite
\usepackage{setspace}
\onehalfspacing
\usepackage{amsmath,amssymb,amsthm}
\usepackage{graphicx,booktabs}
\usepackage{xcolor}
\usepackage{microtype}
% DESIGN: ruhiger Textsatz; Akzentfarbe passend zum CAU-Präsentationsdesign.
\definecolor{akzent}{RGB}{0,135,60}
\definecolor{textgrau}{RGB}{75,85,80}
\addtokomafont{section}{\color{akzent}}
\addtokomafont{subsection}{\color{textgrau}}
\setkomafont{captionlabel}{\sffamily\bfseries\color{textgrau}}
\usepackage[automark,headsepline]{scrlayer-scrpage}
\clearpairofpagestyles
\ihead{\headmark}
\ohead{\textsf{M12}}
\cfoot*{\pagemark}
\setkomafont{pageheadfoot}{\normalfont\sffamily\footnotesize\color{textgrau}}
\pagestyle{scrheadings}
\usepackage{pgfplots}
\pgfplotsset{compat=1.18}
\usepackage[authoryear,round]{natbib}
\newcommand{\quellenfussnote}[1]{\footnote{#1}}
% QUELLEN: Einheitlicher Fußnotenbefehl in Hausarbeit und Präsentation.
% Eine optionale Angabe = Fundstelle; zwei = Vorsatz (z. B. vgl.) + Fundstelle.
\usepackage{xparse}
\NewDocumentCommand{\footcite}{o o m}{%
  \quellenfussnote{%
    \IfNoValueTF{#1}{\citealp{#3}}{%
      \IfNoValueTF{#2}{\citealp[#1]{#3}}{\citealp[#1][#2]{#3}}}%
  }%
}
\usepackage{csquotes}
\usepackage[hidelinks]{hyperref}
\AtBeginDocument{\renewcommand{\refname}{Literaturverzeichnis}}
% Eigene Mathematik-Umgebungen: Nummerierung erfolgt pro Abschnitt.
\newtheorem{satz}{Satz}[section]
\theoremstyle{definition}
\newtheorem{definition}[satz]{Definition}
\newtheorem{beispiel}[satz]{Beispiel}
% BEARBEITEN: wiederverwendbare Angaben; keine Namen im Dokument suchen nötig.
\newcommand{\arbeitstitel}{Quadratische Funktionen im Mathematikunterricht}
\newcommand{\personA}{Vorname Nachname 1}
\newcommand{\personB}{Vorname Nachname 2}
\newcommand{\kurs}{Mathematik · Kurs M12}
\newcommand{\einrichtung}{Justus-Liebig-Universität Gießen}
\newcommand{\betreuung}{Name der betreuenden Lehrkraft}
\newcommand{\abgabedatum}{Abgabedatum}
% LOGO: Datei neben main.tex ablegen und hier den Dateinamen ändern.
% PDF/PNG/JPG möglich; ohne Datei erscheint ein neutraler Platzhalter.
\newcommand{\logodatei}{uni-logo.png}
\newcommand{\hochschullogo}{%
  \IfFileExists{\logodatei}{%
    \includegraphics[width=4cm,height=2cm,keepaspectratio]{\logodatei}%
  }{%
    \begingroup
    \setlength{\fboxsep}{10pt}%
    \fcolorbox{akzent!30}{akzent!4}{%
      \parbox[c][1.3cm][c]{3.2cm}{\centering\sffamily\small
        \textcolor{akzent}{Universitätslogo}\par
        \textcolor{textgrau}{\scriptsize Platzhalter}}}%
    \endgroup
  }%
}
\hypersetup{pdftitle={\arbeitstitel},pdfauthor={\personA, \personB}}

\begin{document}
\begin{titlepage}
  \thispagestyle{empty}
  \singlespacing
  \noindent
  \begin{minipage}[c]{0.60\textwidth}
    {\sffamily\bfseries\large\einrichtung\par}
    \medskip
    {\sffamily\color{textgrau}\kurs\par}
  \end{minipage}\hfill
  \begin{minipage}[c]{0.35\textwidth}
    \raggedleft\hochschullogo
  \end{minipage}

  \vspace{1cm}
  {\color{akzent}\hrule height 1.5pt}
  \vspace{2.1cm}
  {\sffamily\bfseries\color{akzent}\large HAUSARBEIT\par}
  \vspace{0.6cm}
  {\raggedright\sffamily\bfseries\fontsize{29}{35}\selectfont
    \arbeitstitel\par}
  \vspace{0.8cm}
  {\sffamily\large\color{textgrau}Gemeinsam erarbeitet von zwei Personen\par}
  \vfill
  % BEARBEITEN: ggf. Matrikelnummern ergänzen; Vorgaben der Einrichtung prüfen.
  {\renewcommand{\arraystretch}{1.45}
  \begin{tabular}{@{}p{3.4cm}p{11cm}@{}}
    \textsf{Verfasst von} & \personA \\
                         & \personB \\
    \textsf{Kurs}        & \kurs \\
    \textsf{Betreuung}   & \betreuung \\
    \textsf{Abgabe}      & \abgabedatum
  \end{tabular}\par}
  \vspace{1cm}
  {\color{akzent!35}\hrule height 0.6pt}
\end{titlepage}

\pagenumbering{roman}
\section*{Kurzfassung}
% BEARBEITEN: etwa 100--150 Wörter zu Frage, Vorgehen und Ergebnis.
Diese Vorlage zeigt an quadratischen Funktionen grundlegende Elemente
einer wissenschaftlichen Hausarbeit: Fragestellung, mathematische Argumentation,
Unterrichtsplanung und Quellenarbeit. Die Beispiele dienen dem Ausprobieren;
sie sind durch die eigene Untersuchung und deren Ergebnisse zu ersetzen.
\tableofcontents
% OPTIONAL: erst bei mehreren Abbildungen/Tabellen sinnvoll; auskommentieren.
% \listoffigures
% \listoftables
\clearpage
\pagenumbering{arabic}

\section{Einleitung und Fragestellung}
\label{sec:einleitung}
% BEARBEITEN: Motivation, Eingrenzung, Frage und Aufbau der Arbeit erläutern.
Wie können Lernende Nullstellen quadratischer Funktionen rechnerisch bestimmen
und grafisch deuten? Die Beispielstunde verbindet eine gemeinsame Herleitung
an der Tafel mit einer anschließenden Partneraufgabe.

% FORMAT: Leerzeile beginnt einen neuen Absatz; nicht überall \\ verwenden.
\textbf{Fettdruck} hebt sparsam einen Begriff hervor; \emph{Hervorhebung}
betont eine Formulierung. Sonderzeichen schreiben: \%, \&, \_, \#.
Eine Fußnote\footnote{Fußnoten eignen sich für ergänzende Hinweise.
Quellenbelege werden mit dem Befehl \texttt{\textbackslash footcite} gesetzt.}
zeigt eine optionale Ergänzung.

\section{Mathematische Grundlagen}
\label{sec:grundlagen}
% QUELLEN: Die Fußmarke steht unmittelbar an der belegten Aussage.
Spivaks Buch behandelt Grundlagen der Analysis\footcite{spivak2008}. Die folgenden Rechnungen sind eigene
Minimalbeispiele und kein wörtliches Zitat aus diesem Buch.

\begin{definition}[Nullstelle]
  Eine Zahl \(x_0\in\mathbb{R}\) heißt Nullstelle einer Funktion
  \(f\colon\mathbb{R}\to\mathbb{R}\), wenn \(f(x_0)=0\).
\end{definition}

% FORMEL: equation wird nummeriert; equation* oder \[...\] ohne Nummer.
\begin{equation}
  f(x)=ax^2+bx+c, \qquad a\neq 0.
  \label{eq:quadratisch}
\end{equation}
Mit \verb|\eqref{eq:quadratisch}| verweist man automatisch auf
Gleichung~\eqref{eq:quadratisch}. \verb|\label| steht hinter der Nummernquelle,
bei Abbildungen und Tabellen also hinter \verb|\caption|.

\begin{satz}[Nullstellenformel]
  Für \(\Delta=b^2-4ac\geq 0\) liegen die reellen Nullstellen von
  \eqref{eq:quadratisch} bei
  \[
    x_{1,2}=\frac{-b\pm\sqrt{b^2-4ac}}{2a}.
  \]
  Bei \(\Delta=0\) fallen beide Werte zusammen; bei \(\Delta<0\)
  gibt es keine reellen Nullstellen.
\end{satz}
\begin{proof}
  % align richtet Zeilen an & aus. \\ beginnt eine neue Rechenzeile.
  Durch quadratische Ergänzung erhält man
  \begin{align*}
    ax^2+bx+c=0
      &\iff x^2+\frac{b}{a}x=-\frac{c}{a} \\
      &\iff \left(x+\frac{b}{2a}\right)^2
              =\frac{b^2-4ac}{4a^2}.
  \end{align*}
  Für nichtnegative Diskriminante liefert Wurzelziehen die angegebenen Werte
  (das Vorzeichenpaar berücksichtigt auch negatives \(a\)). Für negative
  Diskriminante kann das Quadrat einer reellen Zahl die rechte Seite nicht erreichen.
\end{proof}

\begin{beispiel}
  Für \(f(x)=x^2-4x+3=(x-1)(x-3)\) folgt aus der Nullproduktregel
  \(x_1=1\) und \(x_2=3\). Die Probe ergibt \(f(1)=f(3)=0\).
\end{beispiel}

\subsection{Abbildungen und Tabellen}
% ABBILDUNG: LaTeX darf figure/table verschieben; [htbp] sind Platzierungswünsche.
\begin{figure}[htbp]
  \centering
  % BEARBEITEN: Ausdruck in addplot, Achsen und Beschriftung ändern.
  \begin{tikzpicture}
    \begin{axis}[width=0.75\textwidth,height=6cm,
      axis lines=middle,xlabel={$x$},ylabel={$f(x)$},
      xmin=0,xmax=4,ymin=-1.5,ymax=3.5,grid=major,
      xtick={0,1,2,3,4},ytick={-1,0,1,2,3}]
      \addplot[akzent,thick,domain=0:4,samples=60]{x^2-4*x+3};
      \addplot[only marks,mark=*] coordinates {(1,0) (3,0)};
    \end{axis}
  \end{tikzpicture}
  \caption{Graph von \(f(x)=x^2-4x+3\) mit markierten Nullstellen.
    Quelle: eigene Darstellung.}
  \label{fig:parabel}
\end{figure}
Abbildung~\ref{fig:parabel} ergänzt die Rechnung. Externe Bilder können statt
der TikZ-Grafik mit diesem Befehl eingefügt werden:
\begin{quote}
  \small\verb|\includegraphics[width=0.7\textwidth]{bild.pdf}|
\end{quote}
Die Datei liegt dann neben \texttt{main.tex}; Urheber,
Quelle und gegebenenfalls Lizenz gehören in die Bildunterschrift.

\begin{table}[htbp]
  \centering
  \caption{Wertetabelle zum Beispiel; eigene Berechnung.}
  \label{tab:werte}
  \begin{tabular}{@{}rrrrr@{}}
    \toprule
    \(x\) & 0 & 1 & 2 & 3 \\
    \midrule
    \(f(x)\) & 3 & 0 & -1 & 0 \\
    \bottomrule
  \end{tabular}
\end{table}
Tabelle~\ref{tab:werte} enthält ausgewählte Funktionswerte.

\section{Didaktische Planung der Unterrichtsstunde}
% BEARBEITEN: Klassenstufe, Voraussetzungen, Lernziele, Methoden und Zeiten.
Vorausgesetzt werden der Umgang mit Termen und das Ablesen von Koordinaten.
Als Lernziel bestimmen die Lernenden Nullstellen eines faktorisierbaren
quadratischen Terms und erklären deren Bedeutung am Graphen.
\begin{enumerate}
  \item \textbf{Einstieg (5 Minuten):} Vorwissen zu Nullstellen aktivieren.
  \item \textbf{Erarbeitung (10 Minuten):} Rechenweg gemeinsam an der Tafel entwickeln.
  \item \textbf{Übung (15 Minuten):} Aufgaben in Paaren lösen und kontrollieren.
  \item \textbf{Sicherung (10 Minuten):} Lösungen vergleichen und Fehler besprechen.
  \item \textbf{Abschluss (5 Minuten):} Nullproduktregel in eigenen Worten erklären.
\end{enumerate}
Die Anregung, einen Lösungsweg zu planen und das Ergebnis rückblickend zu prüfen,
lässt sich mit Pólyas Problemlöseheuristik verbinden\footcite{polya1945}.
% Vor Abgabe konkrete Fundstelle in eurer Ausgabe prüfen und ergänzen:
% \footcite[vgl.][S.~23--25]{polya1945} (Seiten hier nur Syntax-Beispiel!).

\subsection{Zusammenarbeit zu zweit}
% BEARBEITEN: tatsächliche Beiträge nennen; beide lesen die gesamte Endfassung.
\begin{itemize}
  \item \personA: Beispielzuständigkeit für Einstieg und Literaturrecherche.
  \item \personB: Beispielzuständigkeit für Herleitung und Aufgabenentwicklung.
  \item Gemeinsam: Unterrichtsplanung, Überarbeitung und Quellenkontrolle.
\end{itemize}

\clearpage
\section{Zitierwerkstatt: Fußnoten in Aktion}
\label{sec:zitierdemo}
\textbf{Layoutdemonstration:} Die folgenden Beispielsätze und Seitenzahlen
zeigen ausschließlich das Zitiersystem. Die Seiten sind nicht als Fundstellen
geprüft; auch der Text in Anführungszeichen ist kein Originalzitat.
Vor der Abgabe diese Demonstration entfernen oder durch belegte Inhalte ersetzen.

\subsection{Einzelne Seiten und Seitenbereiche}
Ein allgemeiner Literaturhinweis braucht keine Seitenangabe: Hier steht ein
Demoverweis auf das gesamte Buch\footcite{polya1945}.

Für eine konkrete Stelle steht die Seite in derselben Fußnote: Diese
Beispielaussage verweist auf eine einzelne Seite\footcite[S.~23]{polya1945}.
Ein zweiter Beleg desselben Buches kann eine andere Fundstelle nennen,
ohne einen zweiten Bibliothekseintrag anzulegen\footcite[S.~42]{polya1945}.

Ein zusammenhängender Abschnitt lässt sich mit einem Seitenbereich
angeben\footcite[S.~23--25]{polya1945}. Dasselbe funktioniert für eine
andere Quelle\footcite[S.~10--12]{spivak2008}.

\subsection{Indirekt, direkt und mehrere Quellen}
Eine sinngemäß wiedergegebene Aussage kann mit einem vorangestellten
\enquote{vgl.} belegt werden\footcite[vgl.][S.~23--25]{polya1945}.
Auch beim zweiten Buch bleibt die Schreibweise gleich\footcite[vgl.][S.~15]{spivak2008}.

\enquote{Dieser erfundene Demotext zeigt nur die Gestaltung eines direkten
Zitats.}\footcite[S.~23]{polya1945}
Ein echtes direktes Zitat benötigt den überprüften Originalwortlaut und die
tatsächliche Fundstelle.

Eine Fußmarke kann mehrere Bücher zusammen belegen\footcite{polya1945,spivak2008}.
Haben die Bücher unterschiedliche Fundstellen, erhält jedes einen eigenen
Beleg\footcite[S.~23]{polya1945}\footcite[S.~10--12]{spivak2008}.
So ist eindeutig, welche Seiten zu welchem Werk gehören.

\subsection{Die Befehle zum Nachbauen}
\begin{quote}
\small
\verb|\footcite{polya1945}|\\
\verb|\footcite[S.~23]{polya1945}|\\
\verb|\footcite[S.~23--25]{polya1945}|\\
\verb|\footcite[vgl.][S.~23--25]{polya1945}|\\
\verb|\footcite{polya1945,spivak2008}|
\end{quote}
Die vollständigen Buchangaben stehen automatisch im Literaturverzeichnis.
Neue Literatur wird in \texttt{literatur.bib} hinterlegt; die genaue Fundstelle
gehört jeweils an die Zitierstelle im Text.

\section{Fazit und Reflexion}
% BEARBEITEN: Frage aus Einleitung beantworten, Grenzen und Ausblick nennen.
Das Beispiel verbindet algebraische Rechnung, Graph und Ergebniskontrolle.
In der eigenen Arbeit ist zu reflektieren, welche Lernziele erreicht wurden
und woran dies erkennbar war. Vor einer tatsächlich gehaltenen Stunde keine
Beobachtungen oder Lernergebnisse als durchgeführt ausgeben.

% LITERATUR: Nur zitierte Einträge erscheinen automatisch. Kein manuelles Tippen.
\clearpage
\addcontentsline{toc}{section}{Literaturverzeichnis}
\bibliographystyle{plainnat}
\bibliography{literatur}

\clearpage
\appendix
\section{Aufgabe und Musterlösung}
% BEARBEITEN: ergänzendes Material; im Haupttext mit \ref{sec:anhang} verweisen.
\label{sec:anhang}
Bestimmt die Nullstellen von \(g(x)=x^2-5x+6\).
\[
  g(x)=(x-2)(x-3)=0 \iff x=2\ \text{oder}\ x=3.
\]
Die Probe ergibt \(g(2)=0\) und \(g(3)=0\).

% OPTIONAL: Erklärung nur verwenden, wenn gefordert, und dann den verbindlichen
% Wortlaut eurer Einrichtung übernehmen (z. B. auch Vorgaben zu KI-Nutzung).
% \clearpage
% \section*{Eigenständigkeitserklärung}
% [Hier den vorgeschriebenen und für euch zutreffenden Wortlaut einsetzen.]
% \vspace{1cm}
% Ort, Datum: \hrulefill
% \vspace{1cm}
% \personA: \hrulefill \qquad \personB: \hrulefill
\end{document}
